% Translated from sample_notes.pdf (handwritten whiteboard photos on bond
% excess returns) by an AI agent. The source was a phone-quality scan of
% marker-on-paper math; the goal is a clean beamer deck that covers the
% same derivation, not a pixel-perfect reproduction of the handwriting.

\documentclass{beamer}

\usetheme{Madrid}
\usepackage{amsmath, amssymb}

\title{Bond Excess Returns and the Expectations Hypothesis}
\author{FBE 633 --- Topic 4 Example}
\date{\today}

\begin{document}

\begin{frame}
  \titlepage
\end{frame}

\begin{frame}{Notation}
  \begin{itemize}
    \item $y^n_t$ = $n$-year continuously compounded yield at date $t$.
    \item $P^n_t = e^{-y^n_t \cdot n}$
      \begin{itemize}
        \item price at time $t$ of an $n$-year zero-coupon bond paying \$1 at time $t+n$.
      \end{itemize}
    \item $r^n_{t+1}$ = log return from holding the $n$-year zero from $t$ to $t+1$:
      \begin{equation*}
        r^n_{t+1} = \log P^{n-1}_{t+1} - \log P^n_t
                  = -(n-1)\, y^{n-1}_{t+1} + n\, y^n_t.
      \end{equation*}
    \item $xr^n_{t+1}$ = log excess return:
      \begin{equation*}
        xr^n_{t+1} = r^n_{t+1} - y^1_t.
      \end{equation*}
  \end{itemize}
\end{frame}

\begin{frame}{Expectations Hypothesis}
  Under EH, log excess returns are unforecastable:
  \begin{equation*}
    xr^n_{t+1} = \alpha + \varepsilon_{t+1}.
  \end{equation*}

  \textbf{Note:}
  \begin{equation*}
    xr^n_{t+1} = -(n-1)\, y^{n-1}_{t+1} + n\, y^n_t - y^1_t.
  \end{equation*}
\end{frame}

\begin{frame}{Implication for Yield Changes}
  EH $\Rightarrow$
  \begin{equation*}
    y^{n-1}_{t+1} - y^n_t
      = -\frac{1}{n-1}\,\alpha \;+\; \frac{y^n_t - y^1_t}{n-1} \;-\; \varepsilon_{t+1}.
  \end{equation*}

  \vspace{0.5em}
  Or, in regression form:
  \begin{equation*}
    y^{n-1}_{t+1} - y^n_t = \alpha_0 + \alpha_1 \cdot \text{slope}_n + e_t,
    \quad \text{with}\ \alpha_1 = 1.
  \end{equation*}

  \vspace{0.5em}
  \begin{center}
    \fbox{\textbf{Higher slope predicts higher future yields.}}
  \end{center}
\end{frame}

\begin{frame}{Why This Matters}
  \begin{itemize}
    \item EH gives a sharp, testable prediction: the slope coefficient is exactly 1.
    \item Campbell-Shiller (1991) and Dai-Singleton (2002) show empirically:
      \begin{itemize}
        \item $\alpha_1$ is \emph{negative}, not $+1$.
        \item It gets \emph{more} negative as maturity $n$ grows.
      \end{itemize}
    \item Interpretation: time-varying risk premia, not just rational expectations of future short rates, drive the yield curve.
    \item We replicate this in section 5 of today's lecture using FRED yields.
  \end{itemize}
\end{frame}

\end{document}
